\section{Timing Analysis of Jam Prevention}
\label{sec:feed}

\noindent\textbf{The feed publishes once a minute and arrives tens of seconds late.} Every response carries a stamp, \texttt{sourceUpdated}, the time its version was published; we polled the feed every 5\,s for 1.9 hours over Mainz and over the area of our drives. Over Mainz, 113 versions arrived 60\,s apart and each became readable 34.4 to 89.5\,s after its stamp; over New Jersey the publication phase shifted, so versions arrived as little as 28\,s apart, each readable 28.9 to 84.0\,s after its stamp. The 123 responses the rig received on the road were 34 to 110\,s old. A vehicle that keys its plan to the feed version (\S~\ref{sec:controller}) applies each version $D=\bar L+P_q$ after its stamp, where $\bar L$ is the largest lag and $P_q$ the query interval: 99.5\,s on the Mainz process. This delivery delay adds to the feed's own probe latency~\cite{wang2018latency}.

\noindent\textbf{The lag condition.} Let $t_v$ be the first time a segment's speed departs from what it would have been without the triggering event, $t_{\mathrm{on}}$ the onset of the jam, and $W=t_{\mathrm{on}}-t_v$ the warning the speed signal gives. Every version stamped before $t_v$ is the same with and without the event.

\begin{proposition}[Lag condition]\label{prop:lag}
With publications at most $g$ seconds apart and each version applied $D$ seconds after its stamp, a feed-only controller changes no advised speed before $t_v+D$, and so changes none before onset if $W<D$. If $D\le W<D+g$ and the stamp phase is uniform, it can act before onset with probability $(W-D)/g$.
\end{proposition}

\begin{proof}[Proof sketch]
The first stamp at or after $t_v$ falls in $[t_v,t_v+g)$ and is applied in $[t_v+D,t_v+D+g)$; nothing applied earlier differs from the case without the event. Acting before onset needs that stamp in $[t_v,t_{\mathrm{on}}-D)$, an interval of length $W-D$ out of $g$.
\end{proof}

The condition is necessary, not sufficient, because a changed advisory still has to change how vehicles drive before the jam forms. Speed makes $W$ short: on the free-flow branch of the speed-density relation, speed hardly changes as density rises until the road breaks down, while density, which a camera measures from the gaps to the vehicles ahead, rises from the start.

\noindent\textbf{How often the local layer must sample.} A rule that observes density only at sample instants protects the critical density only if it acts early enough to absorb the rise that can happen before its next action.

\begin{proposition}[Sampled margin]\label{prop:margin}
Let local density rise no faster than $r$, let a rule sample it every $\Delta$ and act $l$ after a sample, and let its action stop the rise. If the rule acts at the first sample at or above $k^\ast-r(\Delta+l)$, density stays below $k^\ast$.
\end{proposition}

\begin{proof}[Proof sketch]
The sample before the one that triggers the action was below $k^\ast-r(\Delta+l)$, and density rises by at most $r(\Delta+l)$ before the action.
\end{proof}

For the feed, $\Delta+l$ is $60\,\mathrm{s}+D$, about 160\,s on the Mainz process; for a camera sampled every 5\,s with about 1\,s of processing it is 6\,s. The local layer's margin is therefore over 25 times smaller, so it can act close to the critical density.

\noindent\textbf{Why preventing a jam is worth more than draining it.} A jam behaves like a queue at the bottleneck, and the queue's arithmetic decides whether it ever clears.

\begin{proposition}[Point-queue drain]\label{prop:drain}
A queue of $Q$ vehicles with arrival rate $\lambda$ and service rate $\mu$ empties after $Q/(\mu-\lambda)$ if $\mu>\lambda$ and never otherwise. If the bottleneck serves $\mu_f$ before it breaks down and $\mu_q<\mu_f$ after, and $\mu_q\le\lambda<\mu_f$, no queue forms while it has not broken down, and none drains once it has.
\end{proposition}

On the Mainz network the jammed exit serves 3,535\,veh/h under no control, 37\,veh/h less than the background demand of 3,572\,veh/h, so a jam that forms there never drains, and preventing the breakdown is worth more than draining the queue faster.
